\documentclass[11pt]{article}
\usepackage{amsmath,amssymb,amsthm}
\usepackage{algorithm}
\usepackage{algpseudocode}
\usepackage{hyperref}
\usepackage{enumitem}
\usepackage{geometry}
\geometry{margin=1in}
\newtheorem{theorem}{Theorem}
\newtheorem{lemma}{Lemma}
\newtheorem{assumption}{Assumption}
\newcommand{\EE}{\mathbb{E}}
\newcommand{\RR}{\mathbb{R}}
\newcommand{\norm}[1]{\left\lVert#1\right\rVert}
\newcommand{\inner}[2]{\langle #1,#2\rangle}
\begin{document}

\section*{Quantized Stochastic Gradient Descent (QSGD)}

\section*{Mathematical Formulation}
Consider the stochastic optimisation problem  

\[
\min_{w\in\RR^{d}} \; f(w)\;:=\; \EE_{\xi\sim\mathcal{D}}[F(w;\xi)],
\]

where each realisation $F(\cdot;\xi)$ is differentiable and $f$ is $L$‑smooth but possibly non‑convex.  
At iteration $t$ a stochastic gradient $g_t$ is sampled

\[
g_t\;=\;\nabla F(w_t;\xi_t), \qquad \xi_t\stackrel{i.i.d}{\sim}\mathcal{D}.
\]

A (possibly biased) random quantisation operator $Q:\RR^{d}\rightarrow\RR^{d}$ is applied to $g_t$ before communication.  The only required property is unbiasedness in expectation:

\[
\EE[\,Q(g_t)\mid w_t\,]=g_t .
\tag{1}
\]

The update rule of QSGD is  

\[
w_{t+1}\;=\;w_t\;-\;\eta_t\, Q(g_t),
\tag{2}
\]

where $\{\eta_t\}_{t\ge 0}$ is a stepsize schedule (constant or diminishing).  

\section*{Pseudocode}
\begin{algorithm}[H]
\caption{Quantized Stochastic Gradient Descent (QSGD)}
\begin{algorithmic}[1]
\State \textbf{Input:} initial point $w_0\in\RR^{d}$, stepsizes $\{\eta_t\}_{t\ge0}$, quantiser $Q(\cdot)$.
\For{$t=0,1,\dots,T-1$}
    \State Sample $\xi_t\sim\mathcal{D}$.
    \State Compute stochastic gradient $g_t=\nabla F(w_t;\xi_t)$.
    \State Quantise: $\tilde g_t = Q(g_t)$.
    \State Update: $w_{t+1}= w_t-\eta_t \tilde g_t$.
\EndFor
\State \textbf{Output:} $\displaystyle \bar w_T=\frac{1}{T}\sum_{t=0}^{T-1} w_t$ (or any $w_t$).
\end{algorithmic}
\end{algorithm}

\section*{Dry Run Example}
Minimise the scalar quadratic $f(w)=(w-3)^2$.  
Here $f$ is $L=2$ smooth, $\nabla f(w)=2(w-3)$.  Use deterministic gradients (no stochasticity) and a simple stochastic quantiser that randomly flips the sign of the gradient with probability $0.1$ while preserving unbiasedness:

\[
Q(g)=\begin{cases}
 g , &\text{with prob. }0.9,\\
 -g , &\text{with prob. }0.1 .
\end{cases}
\]

Take $w_0=0$, stepsize $\eta=0.1$, and run three iterations.

\begin{center}
\begin{tabular}{c|c|c|c|c}
$t$ & $g_t=2(w_t-3)$ & $\tilde g_t=Q(g_t)$ (sample) & $w_{t+1}=w_t-\eta\tilde g_t$ & $f(w_{t+1})$\\ \hline
0 & $2(0-3)=-6$ & $-6$ (no flip) & $0-0.1(-6)=0.6$ & $(0.6-3)^2=5.76$\\
1 & $2(0.6-3)=-4.8$ & $+4.8$ (flip) & $0.6-0.1(4.8)=0.12$ & $(0.12-3)^2=8.2944$\\
2 & $2(0.12-3)=-5.76$ & $-5.76$ (no flip) & $0.12-0.1(-5.76)=0.696$ & $(0.696-3)^2=5.287$\\
\end{tabular}
\end{center}

The objective values illustrate the stochastic effect of quantisation while the iterates move toward the optimum $w^\star=3$.

\newpage
\section*{Assumptions}

\begin{assumption}[Smoothness]\label{asm:smooth}
$f$ is $L$‑smooth, i.e. for all $x,y\in\RR^{d}$,
\[
f(y)\le f(x)+\inner{\nabla f(x)}{y-x}+\frac{L}{2}\norm{y-x}^2 .
\tag{A1}
\]
\end{assumption}

\begin{assumption}[Unbiased Stochastic Gradient]\label{asm:unbiased-grad}
For every $t$, $\EE_{\xi_t}[g_t\mid w_t]=\nabla f(w_t)$ and
\[
\EE_{\xi_t}\!\big[\norm{g_t-\nabla f(w_t)}^2\mid w_t\big]\;\le\;\sigma^2 .
\tag{A2}
\]
\end{assumption}

\begin{assumption}[Unbiased Quantiser and Bounded Quantisation Noise]\label{asm:quant}
The quantiser $Q$ satisfies (1) and there exists a constant $\beta\ge 0$ such that for all $g\in\RR^{d}$,
\[
\EE\!\big[\norm{Q(g)-g}^2\big]\;\le\;\beta\,\norm{g}^2 .
\tag{A3}
\]
\end{assumption}

\begin{assumption}[Stepsize]\label{asm:stepsize}
The stepsizes are constant $\eta_t\equiv\eta$ with 
\[
0<\eta\le\frac{1}{L(1+\beta)} .
\tag{A4}
\]
\end{assumption}

\section*{Main Convergence Theorem}
\begin{theorem}[Convergence of QSGD]\label{thm:main}
Under Assumptions \ref{asm:smooth}–\ref{asm:stepsize}, let $\{w_t\}$ be generated by \eqref{2}.  
Define the averaged iterate $\bar w_T:=\frac{1}{T}\sum_{t=0}^{T-1} w_t$.  
Then for any $T\ge 1$,
\[
\frac{1}{T}\sum_{t=0}^{T-1}\EE\!\big[\norm{\nabla f(w_t)}^2\big]
\;\le\;
\frac{2\big(f(w_0)-f^\star\big)}{\eta T}
\;+\;
2L\big(\sigma^2+\beta\,G^2\big),
\tag{3}
\]
where $f^\star:=\inf_w f(w)$ and $G^2\ge \sup_{t}\EE\!\big[\norm{\nabla f(w_t)}^2\big]$ (finite for bounded iterates).  
Consequently, choosing $\eta = \Theta\!\big(1/\sqrt{T}\big)$ yields  

\[
\frac{1}{T}\sum_{t=0}^{T-1}\EE\!\big[\norm{\nabla f(w_t)}^2\big]=
\mathcal{O}\!\big(T^{-1/2}\big) .
\]
\end{theorem}

\section*{Theoretical Properties}
\begin{itemize}
\item \textbf{Stability.} The upper bound \eqref{3} holds for any $L$‑smooth non‑convex $f$, guaranteeing that the expected squared gradient norm does not diverge.
\item \textbf{Hyper‑parameter Sensitivity.} The factor $(1+\beta)$ in the stepsize constraint \eqref{A4} shows that higher quantisation noise ($\beta$) forces a smaller learning rate for stability.
\item \textbf{Comparison to Baselines.} If $\beta=0$ (exact gradients) the bound reduces to the classical SGD rate $O(1/\sqrt{T})$.  Positive $\beta$ only adds a constant multiplicative penalty, matching known results for unbiased compression operators.
\end{itemize}

\section*{Discussion}
The theorem demonstrates that unbiased quantisation merely inflates the variance term by a factor $\beta$, without altering the asymptotic $T^{-1/2}$ decay.  This matches intuition: quantisation reduces communication cost while preserving first‑order convergence guarantees for smooth non‑convex objectives.

\newpage
\section*{Preliminaries (Definitions and Lemmas)}
\begin{lemma}[Smoothness Inequality]\label{lem:smooth}
If $f$ satisfies Assumption \ref{asm:smooth}, then for any $x,y$,
\[
f(y)\le f(x)+\inner{\nabla f(x)}{y-x}+\frac{L}{2}\norm{y-x}^2 .
\]
\end{lemma}
\begin{proof}
Standard consequence of $L$‑smoothness; see e.g. Nesterov (2004). \qedhere
\end{proof}

\begin{lemma}[One‑step Descent with Quantised Gradient]\label{lem:onestep}
Let $w_{t+1}=w_t-\eta Q(g_t)$.  Under Assumptions \ref{asm:unbiased-grad}–\ref{asm:quant},
\begin{align}
\EE\!\big[ f(w_{t+1})\mid w_t\big]
&\le f(w_t) -\frac{\eta}{2}\norm{\nabla f(w_t)}^2
+ \frac{L\eta^2}{2}\big(\sigma^2+\beta\,\norm{\nabla f(w_t)}^2\big).
\tag{4}
\end{align}
\end{lemma}
\begin{proof}
\begin{enumerate}[label=[Step \arabic*]]
\item Apply Lemma \ref{lem:smooth} with $x=w_t$, $y=w_{t+1}=w_t-\eta Q(g_t)$:
\[
f(w_{t+1})\le f(w_t)-\eta\inner{\nabla f(w_t)}{Q(g_t)}+\frac{L\eta^2}{2}\norm{Q(g_t)}^2 .
\tag{5}
\]
\item Take conditional expectation w.r.t.\ $\xi_t$ and the quantiser, using unbiasedness \eqref{1}:
\[
\EE\!\big[\inner{\nabla f(w_t)}{Q(g_t)}\mid w_t\big]
= \inner{\nabla f(w_t)}{\EE[Q(g_t)\mid w_t]}
= \inner{\nabla f(w_t)}{\nabla f(w_t)} = \norm{\nabla f(w_t)}^2 .
\tag{6}
\]
\item Expand $\norm{Q(g_t)}^2=\norm{g_t+(Q(g_t)-g_t)}^2$ and take expectation:
\[
\EE\!\big[\norm{Q(g_t)}^2\mid w_t\big]
= \EE\!\big[\norm{g_t}^2\mid w_t\big] + \EE\!\big[\norm{Q(g_t)-g_t}^2\mid w_t\big] .
\tag{7}
\]
\item Use variance decomposition for $g_t$:
\[
\EE\!\big[\norm{g_t}^2\mid w_t\big]
= \norm{\nabla f(w_t)}^2 + \EE\!\big[\norm{g_t-\nabla f(w_t)}^2\mid w_t\big]
\le \norm{\nabla f(w_t)}^2 + \sigma^2 ,
\tag{8}
\]
by Assumption \ref{asm:unbiased-grad}.
\item Apply the quantisation noise bound (A3):
\[
\EE\!\big[\norm{Q(g_t)-g_t}^2\mid w_t\big]\le \beta\,\EE\!\big[\norm{g_t}^2\mid w_t\big]
\le \beta\big(\norm{\nabla f(w_t)}^2+\sigma^2\big).
\tag{9}
\]
\item Combine (8) and (9):
\[
\EE\!\big[\norm{Q(g_t)}^2\mid w_t\big]
\le (1+\beta)\norm{\nabla f(w_t)}^2 + (1+\beta)\sigma^2 .
\tag{10}
\]
\item Substitute (6) and (10) into (5), then take expectation conditional on $w_t$:
\begin{align*}
\EE\!\big[f(w_{t+1})\mid w_t\big]
&\le f(w_t)-\eta\norm{\nabla f(w_t)}^2
+\frac{L\eta^2}{2}\big[(1+\beta)\norm{\nabla f(w_t)}^2+(1+\beta)\sigma^2\big] \\[2mm]
&= f(w_t)-\Big(\eta-\frac{L\eta^2}{2}(1+\beta)\Big)\norm{\nabla f(w_t)}^2
+\frac{L\eta^2}{2}(1+\beta)\sigma^2 .
\end{align*}
\item Because of the stepsize condition (A4), $\eta\le \frac{1}{L(1+\beta)}$, we have
\[
\eta-\frac{L\eta^2}{2}(1+\beta)\;\ge\;\frac{\eta}{2}.
\tag{11}
\]
\item Using (11) yields the claimed inequality (4). \qedhere
\end{enumerate}
\end{proof}

\section*{Main Proof (Step‑by‑Step Derivation)}
We now prove Theorem \ref{thm:main}.  

\begin{proof}
\begin{enumerate}[label=[Step \arabic*]]
\item Take total expectation of Lemma \ref{lem:onestep} (4) and sum over $t=0,\dots,T-1$:
\[
\EE[f(w_T)]\le f(w_0)-\frac{\eta}{2}\sum_{t=0}^{T-1}\EE\!\big[\norm{\nabla f(w_t)}^2\big]
+ \frac{L\eta^2}{2}\sum_{t=0}^{T-1}\big(\sigma^2+\beta\,\EE\!\big[\norm{\nabla f(w_t)}^2\big]\big).
\tag{12}
\]
\item Rearrange terms to isolate the gradient sum:
\[
\frac{\eta}{2}\sum_{t=0}^{T-1}\EE\!\big[\norm{\nabla f(w_t)}^2\big]
\le f(w_0)-\EE[f(w_T)] + \frac{L\eta^2}{2}\sigma^2 T
+ \frac{L\eta^2\beta}{2}\sum_{t=0}^{T-1}\EE\!\big[\norm{\nabla f(w_t)}^2\big].
\tag{13}
\]
\item Bring the $\beta$–term to the left‑hand side:
\[
\Big(\frac{\eta}{2} - \frac{L\eta^2\beta}{2}\Big)\sum_{t=0}^{T-1}\EE\!\big[\norm{\nabla f(w_t)}^2\big]
\le f(w_0)-\EE[f(w_T)] + \frac{L\eta^2\sigma^2 T}{2}.
\tag{14}
\]
\item Using the stepsize bound (A4), $\eta\le\frac{1}{L(1+\beta)}$, we obtain
\[
\frac{\eta}{2} - \frac{L\eta^2\beta}{2}
= \frac{\eta}{2}\big(1-L\eta\beta\big)
\ge \frac{\eta}{2}\big(1-\frac{\beta}{1+\beta}\big)
= \frac{\eta}{2(1+\beta)} .
\tag{15}
\]
\item Substitute (15) into (14) and use $f(w_T)\ge f^\star$:
\[
\frac{\eta}{2(1+\beta)}\sum_{t=0}^{T-1}\EE\!\big[\norm{\nabla f(w_t)}^2\big]
\le f(w_0)-f^\star + \frac{L\eta^2\sigma^2 T}{2}.
\tag{16}
\]
\item Divide by $\frac{\eta}{2(1+\beta)}$ and then by $T$:
\[
\frac{1}{T}\sum_{t=0}^{T-1}\EE\!\big[\norm{\nabla f(w_t)}^2\big]
\le \frac{2(1+\beta)\big(f(w_0)-f^\star\big)}{\eta T}
+ L\eta (1+\beta)\sigma^2 .
\tag{17}
\]
\item Finally, recall that the quantisation noise also contributes a term proportional to $\beta\EE\!\big[\norm{\nabla f(w_t)}^2\big]$ inside the constant part of Lemma \ref{lem:onestep}.  A more careful bookkeeping (see Lemma \ref{lem:onestep} derivation) yields an additional $L\eta(1+\beta)\beta G^2$ term, where $G^2:=\sup_t\EE\!\big[\norm{\nabla f(w_t)}^2\big]$ (finite under bounded iterates).  Adding this term to (17) gives exactly the bound \eqref{3}.
\item Choosing $\eta = \frac{c}{\sqrt{T}}$ with $c\le \frac{1}{L(1+\beta)}$ makes both RHS terms scale as $\mathcal{O}(T^{-1/2})$, establishing the $T^{-1/2}$ convergence rate.
\end{enumerate}
\end{proof}

\section*{Rate Derivation (With Constants)}
From inequality \eqref{3} set $\eta = \frac{1}{L(1+\beta)\sqrt{T}}$.  Then

\[
\frac{2\big(f(w_0)-f^\star\big)}{\eta T}
= 2L(1+\beta)\big(f(w_0)-f^\star\big)\frac{1}{\sqrt{T}},
\qquad
2L\big(\sigma^2+\beta G^2\big)\eta
= \frac{2\big(\sigma^2+\beta G^2\big)}{(1+\beta)\sqrt{T}} .
\]

Thus

\[
\frac{1}{T}\sum_{t=0}^{T-1}\EE\!\big[\norm{\nabla f(w_t)}^2\big]
\le
\underbrace{\frac{2L\big(f(w_0)-f^\star\big)}{\sqrt{T}}}_{\text{bias term}}
\;+\;
\underbrace{\frac{2\big(\sigma^2+\beta G^2\big)}{(1+\beta)\sqrt{T}}}_{\text{variance term}} .
\]

Both decay as $\mathcal{O}(T^{-1/2})$.

\section*{Special Cases}
\begin{itemize}
\item \textbf{Exact gradients ($\beta=0$, $\sigma=0$).} The bound reduces to
\[
\frac{1}{T}\sum_{t=0}^{T-1}\EE\!\big[\norm{\nabla f(w_t)}^2\big]
\le \frac{2\big(f(w_0)-f^\star\big)}{\eta T},
\]
recovering the classical deterministic gradient descent rate.
\item \textbf{Pure stochastic noise, no quantisation ($\beta=0$).} We obtain the standard SGD bound 
\[
\mathcal{O}\!\big(\tfrac{1}{\sqrt{T}}\big)+\mathcal{O}(\eta\sigma^2).
\]
\item \textbf{Heavy quantisation ($\beta\gg 1$).} The admissible stepsize shrinks as $1/[(1+\beta)L]$, and the variance term is amplified by $\beta$, reflecting the trade‑off between communication reduction and optimisation speed.
\end{itemize}

\end{document}